\chapter{Calabi--Yau Manifolds}\label{calabi-yau}
\chapterstatus{complete}

\section{Introduction}\label{calabi-yau:section-introduction}

A compact K\"ahler manifold with vanishing first Chern class carries, in every K\"ahler class, a unique Ricci-flat K\"ahler metric. This is Yau's solution of
the Calabi conjecture (Theorem~\ref{calabi-yau:theorem-yau}). Such manifolds, and in particular those with trivial canonical bundle, are called Calabi--Yau
manifolds. They are among the richest examples in Hodge theory. Their Hodge numbers are constrained but far from trivial, their deformations are unobstructed
and governed by $H^{n-1,1}$, and they are the setting of mirror symmetry. Elliptic curves, K3 surfaces and the quintic threefold are the basic examples. We
compute the Hodge numbers of the quintic threefold completely (Proposition~\ref{calabi-yau:proposition-quintic}). References: \cite{gross-huybrechts-joyce},
\cite[Chapter 6]{joyce-special-holonomy}, \cite{yau-1978}, \cite{beauville-1983}.

\noindent\textbf{Prerequisites.} Chapter~\ref{kodaira} (Vanishing and Embedding Theorems), Chapter~\ref{lefschetz-theorems} (Hard Lefschetz and the
Hodge--Riemann Relations) and Chapter~\ref{characteristic-classes} (Characteristic Classes).

Throughout, $X$ is a compact K\"ahler manifold of dimension $n$ with K\"ahler form $\omega = \frac{i}{2}\sum h_{j\bar k}dz^j \wedge d\bar z^k$ in local holomorphic
coordinates (Proposition~\ref{kahler:proposition-fundamental-form-coordinates}). We write $c_1(X) = c_1(T^{1,0}X) = -c_1(K_X)$.

\section{The Ricci form}\label{calabi-yau:section-ricci}

\begin{definition}\label{calabi-yau:definition-ricci-form}
The \emph{Ricci form} of the K\"ahler metric $\omega$ is the real $(1,1)$-form given in local holomorphic coordinates by
\[
\mathrm{Ric}(\omega) = -i\partial\bar\partial\log\det(h_{j\bar k}) .
\]
The metric is \emph{Ricci-flat} if $\mathrm{Ric}(\omega) = 0$ and \emph{K\"ahler--Einstein} if $\mathrm{Ric}(\omega) = \lambda\omega$ for a constant $\lambda$.
\end{definition}

\begin{proposition}\label{calabi-yau:proposition-ricci-form}
$\mathrm{Ric}(\omega)$ is a well-defined closed real $(1,1)$-form, equal to $2\pi$ times the curvature form $\frac{i}{2\pi}F$ of the anticanonical bundle $K_X^{-1} = \Lambda^nT^{1,0}X$ with
the metric induced by $\omega$. Hence
\[
\left[\frac{1}{2\pi}\mathrm{Ric}(\omega)\right] = c_1(X) \in H^2(X; \RR) .
\]
\end{proposition}

\begin{proof}
In a holomorphic chart, $e = \partial_1 \wedge \cdots \wedge \partial_n$ is a holomorphic frame of $K_X^{-1}$ whose norm for the induced metric is $|e|^2 = \det(g(\partial_j, \bar\partial_k)) = 2^{-n}\det(h_{j\bar k})$.
So the local weight is $\varphi = -\log\det(h_{j\bar k}) + n\log 2$, and by Theorem~\ref{holomorphic-bundles:theorem-chern-connection} the curvature is
$F = \partial\bar\partial\varphi = -\partial\bar\partial\log\det(h_{j\bar k})$, with $c_1(K_X^{-1}) = [\frac{i}{2\pi}F] = [\frac{1}{2\pi}\mathrm{Ric}(\omega)]$. This also shows that
$\mathrm{Ric}(\omega)$ is independent of the chart: under a change of coordinates, $\det(h_{j\bar k})$ is multiplied by $|J|^2$ with $J$ the holomorphic Jacobian determinant, and
$\partial\bar\partial\log|J|^2 = 0$. It is real and closed as a curvature form of a line bundle.
\end{proof}

\begin{remark}\label{calabi-yau:remark-riemannian-ricci}
$\mathrm{Ric}(\omega)(U, V) = \mathrm{Ric}_g(JU, V)$, where $\mathrm{Ric}_g$ is the Ricci tensor of the underlying Riemannian metric \cite[Chapter 4]{huybrechts-complex}. So
Ricci-flat K\"ahler metrics are Ricci-flat Riemannian metrics, and Proposition~\ref{calabi-yau:proposition-ricci-form} shows that $c_1(X)_\RR = 0$ is a necessary
condition for a K\"ahler metric to be Ricci-flat.
\end{remark}

\section{The Calabi conjecture}\label{calabi-yau:section-calabi-conjecture}

\begin{theorem}[Calabi--Yau]\label{calabi-yau:theorem-yau}
Let $X$ be a compact K\"ahler manifold with K\"ahler form $\omega$, and let $\rho$ be a closed real $(1,1)$-form with $[\rho] = 2\pi c_1(X)$. Then there is a unique K\"ahler form
$\omega'$ with $[\omega'] = [\omega]$ and $\mathrm{Ric}(\omega') = \rho$. In particular, if $c_1(X)_\RR = 0$, every K\"ahler class contains a unique Ricci-flat K\"ahler form.
\end{theorem}

Uniqueness is due to Calabi and existence to Yau \cite{yau-1978}. Proposition~\ref{calabi-yau:proposition-monge-ampere} reduces the theorem to a nonlinear elliptic
equation and proves uniqueness; existence is proved by the continuity method with a priori estimates up to order $2$ and $3$ \cite[Chapter 5]{joyce-special-holonomy}.

\begin{proposition}\label{calabi-yau:proposition-monge-ampere}
In the situation of Theorem~\ref{calabi-yau:theorem-yau}, there is a smooth real function $F$ with $\rho = \mathrm{Ric}(\omega) - i\partial\bar\partial F$ and
$\int_X(e^F - 1)\omega^n = 0$. A K\"ahler form $\omega' = \omega + i\partial\bar\partial\varphi$ in the class of $\omega$ satisfies $\mathrm{Ric}(\omega') = \rho$ if and only if
\[
(\omega + i\partial\bar\partial\varphi)^n = e^F\omega^n
\]
(the \emph{complex Monge--Amp\`ere equation}). Solutions of this equation are unique up to additive constants, so $\omega'$ is unique.
\end{proposition}

\begin{proof}
$\rho$ and $\mathrm{Ric}(\omega)$ are cohomologous closed real $(1,1)$-forms, so $\rho = \mathrm{Ric}(\omega) - i\partial\bar\partial F$ by
Corollary~\ref{hodge-decomposition:corollary-ddbar-real}, and adding a constant to $F$ achieves the normalisation. Every K\"ahler form in $[\omega]$ is
$\omega + i\partial\bar\partial\varphi$ for a real $\varphi$, by the same corollary. The ratio of volume forms $\omega'^n/\omega^n = \det(h'_{j\bar k})/\det(h_{j\bar k})$ is a
positive function, so $\mathrm{Ric}(\omega') - \mathrm{Ric}(\omega) = -i\partial\bar\partial\log(\omega'^n/\omega^n)$. Hence $\mathrm{Ric}(\omega') = \rho$ if and only if
$i\partial\bar\partial(\log(\omega'^n/\omega^n) - F) = 0$, if and only if $\log(\omega'^n/\omega^n) = F + c$ with $c$ constant (a real function with $\partial\bar\partial u = 0$
on a compact manifold is pluriharmonic, hence harmonic for the K\"ahler Laplacian, hence constant). Integrating, $\int\omega'^n = \int\omega^n$ because the forms are cohomologous,
and the normalisation of $F$ forces $c = 0$.

For uniqueness, let $\omega_1 = \omega + i\partial\bar\partial\varphi_1$ and $\omega_2 = \omega + i\partial\bar\partial\varphi_2$ with $\omega_1^n = \omega_2^n$, and let $\psi = \varphi_1 - \varphi_2$. Then
\[
0 = \omega_1^n - \omega_2^n = i\partial\bar\partial\psi \wedge T, \qquad T = \sum_{k=0}^{n-1}\omega_1^k \wedge \omega_2^{n-1-k} .
\]
Multiply by $\psi$ and integrate. Since $T$ is closed, Stokes' theorem gives
$0 = \int_X\psi\,i\partial\bar\partial\psi \wedge T = -\int_Xi\,\partial\psi \wedge \bar\partial\psi \wedge T$. At each point, $\omega_1$ and $\omega_2$ are positive, so $T$ is a sum of
products of positive $(1,1)$-forms, and $i\,\partial\psi \wedge \bar\partial\psi \wedge T$ is a non-negative multiple of the volume form, positive where $\partial\psi \neq 0$ (diagonalise
$\omega_1$ and $\omega_2$ simultaneously at the point). Hence $\partial\psi = 0$, so $d\psi = 0$ as $\psi$ is real, and $\psi$ is constant.
\end{proof}

\section{Calabi--Yau manifolds and their Hodge numbers}\label{calabi-yau:section-definition}

\begin{definition}\label{calabi-yau:definition-calabi-yau}
A \emph{Calabi--Yau manifold} is a compact K\"ahler manifold $X$ with trivial canonical bundle $K_X \cong \mathcal{O}_X$. Equivalently, there is a nowhere vanishing holomorphic
$n$-form $\Omega$, the \emph{holomorphic volume form}. It is \emph{strict} if moreover $h^{p,0}(X) = 0$ for $0 < p < n$.
\end{definition}

\begin{proposition}\label{calabi-yau:proposition-basic}
Let $X$ be a Calabi--Yau manifold of dimension $n$.
\begin{enumerate}
\item $c_1(X) = 0$, so $X$ carries Ricci-flat K\"ahler metrics (Theorem~\ref{calabi-yau:theorem-yau}). Moreover $h^{0,q} = h^{n,q}$ for all $q$; in particular
$h^{n,0} = h^{0,n} = 1$.
\item Contraction with $\Omega$ is an isomorphism of holomorphic bundles $T^{1,0}X \to \Omega^{n-1}_X$, $v \mapsto \iota_v\Omega$, and more generally
$\Lambda^pT^{1,0}X \cong \Omega^{n-p}_X$. Hence $H^1(X, T_X) \cong H^{n-1,1}(X)$.
\end{enumerate}
\end{proposition}

\begin{proof}
(1) $c_1(X) = -c_1(K_X) = 0$, and $h^{n,q} = h^q(X, K_X) = h^q(X, \mathcal{O}_X) = h^{0,q}$ (Theorem~\ref{hodge-decomposition:theorem-hodge-decomposition}(4)). For $q = 0$ this is
$1$, and $h^{0,n} = h^{n,0}$ by Hodge symmetry. (2) At a point choose coordinates with $\Omega = f\,dz^1 \wedge \cdots \wedge dz^n$, $f \neq 0$. Then $\iota_{\partial_j}\Omega = (-1)^{j-1}f\,dz^1 \wedge \cdots \widehat{dz^j} \cdots \wedge dz^n$,
and these form a basis of $\Omega^{n-1}$, so contraction is a holomorphic bundle isomorphism; the same computation with $\iota_{\partial_{j_1}} \cdots \iota_{\partial_{j_p}}$
gives $\Lambda^pT^{1,0}X \cong \Omega^{n-p}_X$. By Dolbeault's theorem, $H^1(X, T_X) \cong H^1(X, \Omega^{n-1}) \cong H^{n-1,1}(X)$.
\end{proof}

\begin{example}\label{calabi-yau:example-calabi-yau}
\begin{enumerate}
\item Complex tori $\CC^n/\Lambda$, with $\Omega = dz^1 \wedge \cdots \wedge dz^n$ and the flat metric, which is Ricci-flat. They are not strict for $n \geq 2$.
\item Elliptic curves are the Calabi--Yau manifolds of dimension $1$, since $K$ is trivial exactly in genus $1$.
\item A \emph{K3 surface} is a strict Calabi--Yau surface, that is, $K_X \cong \mathcal{O}$ and $h^{1,0} = 0$. Smooth quartic surfaces in $\CC\PP^3$ are K3 surfaces
(Example~\ref{hodge-decomposition:example-quartic}).
\item A smooth hypersurface of degree $n + 2$ in $\CC\PP^{n+1}$ has trivial canonical bundle by adjunction, $K = \mathcal{O}(n + 2 - n - 2) = \mathcal{O}$
(Proposition~\ref{holomorphic-bundles:proposition-adjunction}), and for $n \geq 2$ it is strict by the Lefschetz hyperplane theorem
(Theorem~\ref{kodaira:theorem-lefschetz-hyperplane}), since $h^{p,0}(X) = h^{p,0}(\CC\PP^{n+1}) = 0$ for $0 < p < n$.
\item Products of Calabi--Yau manifolds are Calabi--Yau, with $\Omega = \mathrm{pr}_1^*\Omega_1 \wedge \mathrm{pr}_2^*\Omega_2$; they are not strict.
\end{enumerate}
\end{example}

\begin{proposition}[The quintic threefold]\label{calabi-yau:proposition-quintic}
Let $X \subseteq \CC\PP^4$ be a smooth quintic hypersurface. Then $X$ is a strict Calabi--Yau threefold with Hodge diamond determined by
\[
h^{0,0} = h^{3,0} = 1, \qquad h^{1,0} = h^{2,0} = 0, \qquad h^{1,1} = 1, \qquad h^{2,1} = 101,
\]
and Euler characteristic $\chi(X) = -200$.
\end{proposition}

\begin{proof}
$K_X$ is trivial and $h^{1,0} = h^{2,0} = 0$ by Example~\ref{calabi-yau:example-calabi-yau}(4), with $n = 3$. By the Lefschetz hyperplane theorem, $b_1(X) = 0$ and $b_2(X) = b_2(\CC\PP^4) = 1$,
so $h^{1,1} = b_2 - 2h^{2,0} = 1$. By Poincar\'e duality $b_4 = 1$ and $b_5 = 0$, and by the Hodge decomposition $b_3 = h^{3,0} + h^{2,1} + h^{1,2} + h^{0,3} = 2 + 2h^{2,1}$. So
$\chi(X) = 1 + 1 - b_3 + 1 + 1 = 2 - 2h^{2,1}$.

We compute $\chi(X) = \int_Xc_3(TX)$ (Theorem~\ref{characteristic-classes:theorem-euler}). Let $h$ be the hyperplane class. The tangent bundle of $\CC\PP^4$ has total Chern class $(1 + h)^5$
(Example~\ref{characteristic-classes:example-tautological}). The normal bundle of $X$ is $\mathcal{O}(X)|_X = \mathcal{O}(5)|_X$ (Proposition~\ref{holomorphic-bundles:proposition-adjunction}),
with $c = 1 + 5h$. From $0 \to TX \to T\CC\PP^4|_X \to N \to 0$ and the Whitney formula,
\[
c(TX) = \frac{(1 + h)^5}{1 + 5h} = (1 + 5h + 10h^2 + 10h^3)(1 - 5h + 25h^2 - 125h^3) = 1 + 10h^2 - 40h^3
\]
restricted to $X$, modulo $h^4$. So $c_1(X) = 0$, as it must be, and $\chi(X) = -40\int_Xh^3$. Finally $\int_Xh^3 = 5$: the class of $X$ is Poincar\'e dual to $c_1(\mathcal{O}(X)) = 5h$
\cite[Chapter 1, Section 1]{griffiths-harris}, and $\int_{\CC\PP^4}h^4 = 1$. Hence $\chi(X) = -200$ and $h^{2,1} = 101$.
\end{proof}

\begin{remark}\label{calabi-yau:remark-moduli-count}
The number $101$ is the number of moduli of quintic threefolds: the quintic polynomials in five variables form a space of dimension $\binom{9}{4} = 126$, and removing scalars and the
$24$-dimensional group $\mathrm{PGL}_5$ leaves $126 - 1 - 24 = 101$ parameters. This matches $h^1(X, T_X) = h^{2,1}$ (Proposition~\ref{calabi-yau:proposition-basic}), which is the
dimension of the space of first-order deformations of the complex structure. By the Bogomolov--Tian--Todorov theorem the deformations of a Calabi--Yau manifold are
unobstructed, so its Kuranishi space is smooth of dimension $h^{n-1,1}$ \cite{tian-1987}, \cite{todorov-1989}.
\end{remark}

\begin{proposition}[Complete intersection Calabi--Yau threefolds]\label{calabi-yau:proposition-complete-intersections}
Let $X = V(f_1, \ldots, f_r) \subseteq \CC\PP^{3+r}$ with $f_i$ homogeneous of degree $d_i \geq 2$ and $\sum_id_i = 4 + r$, and assume that every partial intersection $V(f_1, \ldots, f_j)$, $1 \leq j \leq r$, is smooth of
codimension $j$. Then $X$ is a strict Calabi--Yau threefold with $h^{1,1} = 1$, and
\[
\begin{array}{c|ccccc}
(d_1, \ldots, d_r) & (5) & (2, 4) & (3, 3) & (2, 2, 3) & (2, 2, 2, 2) \\ \hline
c_2(X) & 10h^2 & 7h^2 & 6h^2 & 5h^2 & 4h^2 \\
c_3(X) & -40h^3 & -22h^3 & -16h^3 & -12h^3 & -8h^3 \\
\int_Xh^3 = \prod_id_i & 5 & 8 & 9 & 12 & 16 \\
\chi(X) & -200 & -176 & -144 & -144 & -128 \\
h^{2,1}(X) & 101 & 89 & 73 & 73 & 65
\end{array}
\]
These five are all the possibilities with every $d_i \geq 2$.
\end{proposition}

\begin{proof}
Write $X_j = V(f_1, \ldots, f_j)$, so that $X_j \subseteq X_{j-1}$ is a smooth hypersurface cut out by a section of $\mathcal{O}(d_j)$. By adjunction (Proposition~\ref{holomorphic-bundles:proposition-adjunction}) applied
$r$ times, $K_X = \mathcal{O}(\sum_id_i - 4 - r)|_X = \mathcal{O}_X$. By the Lefschetz hyperplane theorem (Theorem~\ref{kodaira:theorem-lefschetz-hyperplane}) applied to each $X_j \subseteq X_{j-1}$, $H^k(X) \cong H^k(\CC\PP^{3+r})$
for $k < 3$ (with $\QQ$-coefficients); so $b_1 = 0$ and $b_2 = 1 = 2h^{2,0} + h^{1,1}$, and since $h^{1,1} \geq 1$ (the K\"ahler class is nonzero,
Proposition~\ref{kahler:proposition-kahler-class-powers}), $h^{1,0} = h^{2,0} = 0$ and $h^{1,1} = 1$, which makes $X$ strict. As in the proof of Proposition~\ref{calabi-yau:proposition-quintic}, $b_4 = 1$, $b_3 = 2 + 2h^{2,1}$ and
$\chi(X) = 2 - 2h^{2,1}$.

The normal bundle of $X$ is $\bigoplus_i\mathcal{O}(d_i)|_X$ (adjunction at each step), so by the Whitney formula $c(TX) = (1 + h)^{4+r}\prod_i(1 + d_ih)^{-1}$ restricted to $X$, modulo $h^4$. For instance for $(2, 4)$,
\[
(1 + 6h + 15h^2 + 20h^3)(1 - 6h + 28h^2 - 120h^3) = 1 + 7h^2 - 22h^3,
\]
using $(1 + 2h)^{-1}(1 + 4h)^{-1} = (1 - 2h + 4h^2 - 8h^3)(1 - 4h + 16h^2 - 64h^3)$ modulo $h^4$; the other columns are the same computation. The degree $\int_Xh^3 = \prod_id_i$ holds because the class of $X$ is
$\prod_id_i\,h^r$, each hypersurface having class $d_ih$ (as in the proof of Proposition~\ref{calabi-yau:proposition-quintic}). Then $\chi(X) = \int_Xc_3(TX)$ (Theorem~\ref{characteristic-classes:theorem-euler}) gives
the fourth row, and $h^{2,1} = 1 - \chi/2$ the last. Finally, $\sum_i(d_i - 1) = 4$ with all $d_i \geq 2$ has exactly the five solutions listed, up to order.
\end{proof}

\section{Structure and symmetries}\label{calabi-yau:section-structure}

\begin{theorem}[Beauville--Bogomolov decomposition]\label{calabi-yau:theorem-beauville-bogomolov}
Let $X$ be a compact K\"ahler manifold with $c_1(X)_\RR = 0$. Then there is a finite \'etale cover $\tilde X \to X$ with
\[
\tilde X \cong T \times \prod_iY_i \times \prod_jZ_j,
\]
where $T$ is a complex torus, the $Y_i$ are simply connected strict Calabi--Yau manifolds with holonomy $\mathrm{SU}(\dim Y_i)$, and the $Z_j$ are simply connected
\emph{hyperk\"ahler} manifolds, with holonomy $\mathrm{Sp}(\dim Z_j/2)$ and $H^0(Z_j, \Omega^2)$ spanned by a symplectic form.
\end{theorem}

This is \cite{beauville-1983}. Its proof uses the Ricci-flat metric of Theorem~\ref{calabi-yau:theorem-yau} and the de Rham and Berger classifications of holonomy
groups. In dimension $2$ the only simply connected factors are K3 surfaces, which are both strict Calabi--Yau and hyperk\"ahler.

\begin{remark}\label{calabi-yau:remark-bochner}
The first step is Bochner's argument. For a Ricci-flat metric, harmonic $1$-forms are parallel (Theorem~\ref{hodge-theorem:theorem-bochner}), and more generally every holomorphic
$p$-form is parallel. So the holomorphic forms are determined by the holonomy representation, and $h^{p,0}$ is the dimension of the space of holonomy-invariant forms of type
$(p, 0)$ at a point \cite{gross-huybrechts-joyce}. For holonomy exactly $\mathrm{SU}(n)$ the only invariants are $1$ and the volume form, which gives $h^{p,0} = 0$ for $0 < p < n$.
\end{remark}

\begin{remark}\label{calabi-yau:remark-mirror-symmetry}
\emph{Mirror symmetry} predicts that strict Calabi--Yau threefolds come in pairs $(X, X^\vee)$ with $h^{1,1}(X) = h^{2,1}(X^\vee)$ and $h^{2,1}(X) = h^{1,1}(X^\vee)$, the
complex geometry of $X$ being equivalent to the symplectic geometry of $X^\vee$. The mirror of the quintic has $h^{1,1} = 101$ and $h^{2,1} = 1$, and the resulting prediction of the
numbers of rational curves on the quintic \cite{candelas-1991} was a starting point of enumerative mirror symmetry.
\end{remark}

\begin{remark}\label{calabi-yau:remark-hodge-conjecture}
For projective Calabi--Yau manifolds the Hodge conjecture is known in some degrees. On any projective threefold it holds in all degrees, by the Lefschetz $(1,1)$-theorem
in degree $2$ and hard Lefschetz in degree $4$ (Chapter~\ref{lefschetz-11}); the interesting group $H^3$ carries no Hodge classes, being of odd degree. For K3 surfaces it holds
in all degrees (Chapter~\ref{known-cases}). For higher-dimensional hyperk\"ahler varieties it is largely open.
\end{remark}

\section{Exercises}\label{calabi-yau:section-exercises}

\begin{exercise}\label{calabi-yau:exercise-k3-hodge}
Show that a smooth quartic surface $X \subseteq \CC\PP^3$ has $\chi(X) = 24$ by computing $c_2(TX)$ as in Proposition~\ref{calabi-yau:proposition-quintic}, and deduce $h^{1,1} = 20$.
\end{exercise}

\begin{solution}
$c(TX) = (1 + h)^4/(1 + 4h) = (1 + 4h + 6h^2)(1 - 4h + 16h^2) = 1 + 6h^2$ modulo $h^3$, so $c_1 = 0$ and $\chi(X) = 6\int_Xh^2 = 6 \cdot 4 = 24$. With $b_1 = b_3 = 0$ and $b_0 = b_4 = 1$
we get $b_2 = 22$, and $h^{1,1} = b_2 - 2h^{2,0} = 22 - 2 = 20$.
\end{solution}

\begin{exercise}\label{calabi-yau:exercise-ricci-projective}
Show that the Fubini--Study metric on $\CC\PP^n$ is K\"ahler--Einstein with $\mathrm{Ric}(\omega_{FS}) = 2\pi(n + 1)\omega_{FS}$.
\end{exercise}

\begin{solution}
By $U(n+1)$-invariance (Proposition~\ref{kahler:proposition-fubini-study}), $\mathrm{Ric}(\omega_{FS})$ is an invariant closed real $(1,1)$-form. At the point $w = 0$ of the affine
chart, $\omega_{FS} = \frac{i}{2\pi}\partial\bar\partial\log(1 + |w|^2)$ has $h_{j\bar k} = \frac1\pi\frac{(1 + |w|^2)\delta_{jk} - \bar w_jw_k}{(1 + |w|^2)^2}$, whose determinant is
$\pi^{-n}(1 + |w|^2)^{-n-1}$. So $\mathrm{Ric}(\omega_{FS}) = (n + 1)i\partial\bar\partial\log(1 + |w|^2) = 2\pi(n + 1)\omega_{FS}$ on the chart, hence everywhere.
\end{solution}

\begin{exercise}\label{calabi-yau:exercise-hypersurface-euler}
Show that a smooth hypersurface $X \subseteq \CC\PP^{n+1}$ of degree $d$ has $c_1(X) = (n + 2 - d)h|_X$, so that it is Fano ($c_1 > 0$), Calabi--Yau or of general type
($c_1 < 0$) according as $d < n + 2$, $d = n + 2$ or $d > n + 2$.
\end{exercise}

\begin{solution}
$c(TX) = (1 + h)^{n+2}/(1 + dh)$ has degree-one part $(n + 2)h - dh$. Alternatively, $K_X = \mathcal{O}(d - n - 2)|_X$ by adjunction.
\end{solution}
